\def\sode{4}
\begin{name}
	{\tenchude}
	{\tendethi}
	{\tentruong}
	{\thoigian}
\end{name}
\Opensolutionfile{ans}[ans/ans-HXN-\sode-T]
\caulc
\begin{ex}%Câu 1
	Cho hàm số $y=f(x)$ xác định trên $\mathbb{R}$ và có bảng biến thiên như hình vẽ sau:
	\begin{center}
		\begin{tikzpicture}[>=stealth]
			\tkzTabInit[nocadre=false,lgt=1,espcl=2.5,deltacl=0.5]{$x$/.7 ,$y'$/.7,$y$/2}
			{$-\infty$ , $-1$ , $2$ , $+\infty$}
			\tkzTabLine{ , - , $0$ , + , $0$ , - , }
			\tkzTabVar{+/$+\infty$ , -/$-2$ , +/$4$ , -/$-\infty$}
		\end{tikzpicture}
	\end{center}
	Hàm số $y=f(x)$ đồng biến trên khoảng nào?
	\choice
	{$\left(-\infty\,;\,\,-1\right)$}
	{$\left(-2\,;\,\,4\right)$}
	{$\left(2\,;\,\,+\infty\right)$}
	{\True $\left(-1\,;\,\,2\right)$}
	\loigiai{
		Chọn D.}
\end{ex}
\begin{ex}%Câu 2
	Trong không gian $Oxyz$ , cho điểm $ M\left(2\,;\,\,-1\,;\,\,4\right)$ và mặt phẳng $(P):\,3x-2y+z+1=0$. Phương trình của mặt phẳng đi qua $ M$ và song song với mặt phẳng $(P)$ là
	\choice
	{$ 2x-2y+4z-21=0$}
	{$ x-2z+1=0$}
	{$10x+9y+5z-74=0$}
	{\True $ 3x-2y+z-12=0$}
	\loigiai{
		Chọn D.\\
		Phương trình của mặt phẳng đi qua $ M$ và song song với mặt phẳng $(P)$ là\\
		$ 3\left(x-2\right)-2\left(y+1\right)+\left(z-4\right)=0$$\Leftrightarrow 3x-2y+z-12=0$.}
\end{ex}
\begin{ex}%Câu 3
	Điều tra về mức lương khởi điểm (đơn vị: triệu đồng) của $ 20$ công nhân, ta có bảng số liệu sau\\
	\centerline{\begin{tabular}{|c|c|c|c|c|c|}
			\hline
			Mức lương & $\left[5\,;\,\,6\right)$ & $\left[6\,;\,\,7\right)$ & $\left[7\,;\,\,8\right)$ & $\left[8\,;\,\,9\right)$ & $\left[9\,;\,\,10\right)$ \\
			\hline
			Tần số    & $ 4$                     & $ 5$                     & $ 5$                     & $ 4$                     & $ 2$                      \\
			\hline
		\end{tabular}}\\
	Phương sai của mẫu số liệu ghép nhóm là (làm tròn đến hàng phần trăm):
	\choice
	{$s^2=0,63$}
	{$s^2=2,52$}
	{$s^2=1,26$}
	{\True $s^2=1,59$}
	\loigiai{
		Chọn D.\\
		Ta có bảng thống kê về mức lương theo giá trị đại diện như sau:\\
		\centerline{\begin{tabular}{|c|c|c|c|c|c|}
				\hline
				Mức lương        & $\left[5\,;\,\,6\right)$ & $\left[6\,;\,\,7\right)$ & $\left[7\,;\,\,8\right)$ & $\left[8\,;\,\,9\right)$ & $\left[9\,;\,\,10\right)$ \\
				\hline
				Giá trị đại diện & $ 5,5$                   & $ 6,5$                   & $ 7,5$                   & $ 8,5$                   & $ 9,5$                    \\
				\hline
				Tần số           & $ 4$                     & $ 5$                     & $ 5$                     & $ 4$                     & $ 2$                      \\
				\hline
			\end{tabular}}\\
		Số trung bình của mẫu số liệu ghép nhóm là\\
		$\bar{x}=\dfrac{4\cdot 5,5+5\cdot 6,5+5\cdot 7,5+4\cdot 8,5+2\cdot 9,5}{20}=7,25$.\\
		Phương sai của mẫu số liệu ghép nhóm là $$s^2=\dfrac{4\cdot{\left(5,5-7,25\right)^2}+5\cdot{\left(6,5-7,25\right)^2}+5\cdot{\left(7,5-7,25\right)^2}+4\cdot{\left(8,5-7,25\right)^2}+2\cdot{\left(9,5-7,25\right)^2}}{20}\approx 1,59$$.}
\end{ex}
\begin{ex}%Câu 4
	Hàm số nào sau đây không là một nguyên hàm của hàm số $f(x)=3x^2+2x-1$ .
	\choice
	{\True $F(x)=x^3+x^2-1$}
	{$F(x)=x^3+x^2-x$}
	{$F(x)=x^3+x^2-x+2025$}
	{$F(x)=x^3+x^2-x-1$}
	\loigiai{
		Chọn A.\\
		$F(x)=\displaystyle\int{f(x)\text{d}x}=\displaystyle\int{\left(3x^2+2x-1\right)\text{d}x}=x^3+x^2-x+C$ (Với $C$ là một hằng số).}
\end{ex}
\begin{ex}%Câu 5
	Trong không gian $ Oxyz$cho ba điểm $ M\left(1\,;\,\,1\,;\,\,1\right),\,\,N\left(2\,;\,\,3\,;\,\,4\right),\,\,P\left(7\,;\,\,7\,;\,\,5\right)$. Tìm tọa độ điểm $ Q$ để tứ giác $ MNPQ$ là hình bình hành
	\choice
	{\True $ Q\left(6\,;\,\,5\,;\,\,2\right)$}
	{$ Q\left(-6\,;\,\,-5\,;\,\,-2\right)$}
	{$ Q\left(-2\,;\,\,-3\,;\,\,-4\right)$}
	{$Q\left(-4\,;\,\,-3\,;\,\,0\right)$}
	\loigiai{
		Chọn A.\\
		Ta có $\overrightarrow{MN}=\left(1\,;\,\,2\,;\,\,3\right)\,,\,\,\overrightarrow{QP}=\left(7-x_Q\,;\,\,7-y_Q\,;\,\,5-z_Q\right)$.\\
		$ MNPQ$ là hình bình hành $\Leftrightarrow\overrightarrow{MN}=\overrightarrow{QP}$$\Leftrightarrow\left\{\begin{aligned}
			 & 1=7-x_Q \\
			 & 2=7-y_Q \\
			 & 3=5-z_Q \\
		\end{aligned}\right.\Leftrightarrow\left\{\begin{aligned}
			 & {x_Q}=6 \\
			 & {y_Q}=5 \\
			 & {z_Q}=2 \\
		\end{aligned}\right.$. Vậy $ Q(6;5;2)$.}
\end{ex}
\begin{ex}%Câu 6
	Cho cấp số cộng $\left(u_n\right)$ có số hạng đầu $u_1=\dfrac{1}{4}$ và công sai $d=-\dfrac{1}{4}$ . Tổng 5 số hạng đầu tiên của cấp số cộng là
	\choice
	{$S_5=\dfrac{5}{4}$}
	{$S_5=\dfrac{4}{5}$}
	{\True $S_5=-\dfrac{5}{4}$}
	{$S_5=-\dfrac{4}{5}$}
	\loigiai{
		Chọn C.\\
		Sử dụng công thức $S_n=\dfrac{n\left[2u_1+\left(n-1\right)d\right]}{2}$ , ta có: $S_5=\dfrac{5\left[2\cdot\dfrac{1}{4}+\left(5-1\right)\cdot\left(-\dfrac{1}{4}\right)\right]}{2}=-\dfrac{5}{4}$ .}
\end{ex}
\begin{ex}%Câu 7
	Nghiệm của phương trình $3^x=10$ là
	\choice
	{$ x=\dfrac{10}{3}$}
	{\True $ x=\log_310$}
	{$ x=\log_{10}3$}
	{$ x=\dfrac{10}{3}$}
	\loigiai{
		Chọn B.\\
		Ta có $3^x=10\Leftrightarrow x=\log_310$.}
\end{ex}
\begin{ex}%Câu 8
	Cho hình chóp $ S.ABCD$ có đáy $ ABCD$ là hình vuông và $ SA\perp\left(ABCD\right)$. Mặt phẳng $\left(SBC\right)$ vuông góc với mặt phẳng nào sau đây?
	\choice
	{$\left(SCD\right)$}
	{$\left(ABCD\right)$}
	{\True $\left(SAB\right)$}
	{$\left(SBD\right)$}
	\loigiai{
		Chọn C.\\
		Ta có $\left\{\begin{aligned}
				 & BC\perp AB \\
				 & BC\perp SA \\
			\end{aligned}\right.$nên $ BC\perp\left(SAB\right)$;\\
		mà $ BC\subset\left(SBC\right)\Rightarrow (SAB) \perp (SBC) $.}
\end{ex}
\begin{ex}%Câu 9
	Đường tiệm cận xiên của đồ thị hàm số $ y=\dfrac{x^2-3x+4}{x+2}$ là đường thẳng
	\choice
	{$ y=-x+1$}
	{$ y=x-1$}
	{\True $ y=x-5$}
	{$ y=-x-5$}
	\loigiai{
		Chọn C.\\
		Ta có $ y=\dfrac{x^2-3x+4}{x+2}$$=x-5+\dfrac{14}{x+2}$; mặt khác $\underset{x\to+\infty}{\lim}\,\left[y-\left(x-5\right)\right]=\underset{x\to+\infty}{\lim}\,\dfrac{14}{x+2}=0$.\\
			Do đó tiệm cận xiên của đồ thị hàm số là đường thẳng $y=x-5 $.}
\end{ex}
\begin{ex}%Câu 10
	Cho tứ diện $ABCD$ . Gọi $M\,,\,\,N$ lần lượt là trung điểm của $AD$ và $BC$ . Tổng $\overrightarrow{AB}+\overrightarrow{DC}$ bằng
	\choice
	{$\vec{0}$}
	{$ 2\overrightarrow{AD}$}
	{$ 2\overrightarrow{NM}$}
	{\True $ 2\overrightarrow{MN}$}
	\loigiai{
		Chọn D.\\
		Ta có: $\overrightarrow{AB}+\overrightarrow{DC}=\overrightarrow{AM}+\overrightarrow{MN}+\overrightarrow{NB}+\overrightarrow{DM}+\overrightarrow{MN}+\overrightarrow{NC}$ $=\left(\overrightarrow{AM}+\overrightarrow{DM}\right)+2\overrightarrow{MN}+\left(\overrightarrow{NB}+\overrightarrow{NC}\right)=2\vec{MN}.$\\
		(Vì $M,\,N$ lần lượt là trung điểm của $AD$ và $BC$ nên $\overrightarrow{AM}+\overrightarrow{DM}=\overrightarrow{0}\,,\,\,\overrightarrow{NB}+\overrightarrow{NC}=\overrightarrow{0}$).}
\end{ex}
\begin{ex}%Câu 11
	Tiếp tuyến của đồ thị hàm số$ y=x^3-3x^2-2$ có hệ số góc $ k=-3$ có phương trình là
	\choice
	{$ y=-3x-7$}
	{$ y=-3x+7$}
	{$ y=-3x+1$}
	{\True $ y=-3x-1$}
	\loigiai{
		Chọn D\\
		Đạo hàm $y'=3x^2-6x$. Gọi $\left(x_0\,;\,\,y_0\right)$ là tiếp điểm của tiếp tuyến với đồ thị hàm số.\\
		Hệ số góc tiếp tuyến $ k=-3$ nên $ 3x_0^2-6x_0=-3\Rightarrow x_0^2-2x_0+1=0\Rightarrow{x_0}=1\Rightarrow{y_0}=-4$.\\
		Phương trình tiếp tuyến của đồ thị hàm số là: $ y=-3\left(x-1\right)-4\,\,\,\text{hay} y=-3x-1$$ $.}
\end{ex}
\begin{ex}%Câu 12
	\immini[thm]{ Đường gấp khúc trong hình vẽ là đồ thì hàm số $ y=f(x)$ trên đoạn $\left[-2\,;\,\,3\right]$. Tích phân $\displaystyle\int\limits_{-2}^3f(x)\text{d}x$ bằng
		\choice
		{\True $\dfrac{13}{2}$}
		{$\dfrac{17}{2}$}
		{$\dfrac{15}{2}$}
		{$\dfrac{5}{2}$}}{\includegraphics[scale=.8]{img/HXN-4.12}}
	\loigiai{
	Chọn A.\\
	Giá trị tích phân $\displaystyle\int\limits_{-2}^3f(x)\text{d}x$ chính là tổng diện tích hai tam giác ABC và CDE như hình vẽ.\\
	Ta có: $\displaystyle\int\limits_{-2}^3f(x)\text{d}x=S_{ABK}+S_{BCL}=\dfrac{1}{2}\cdot 3\cdot 3+\dfrac{1}{2}\cdot 2\cdot 2=\dfrac{13}{2}$.}
\end{ex}

\Closesolutionfile{ans}
\cauds
\Opensolutionfile{ans}[ans/ans-HXN-\sode-TF]


\begin{ex}
	\immini[thm]{Một mô hình trò chơi vòng quay ở công viên có chiều cao tối đa 23 $m$ so với mặt đất, bán kính vòng quay là 20 $m$. Hai bạn Hoa và Mai cùng chơi chung lượt quay và ngồi trong hai cabin B, C mà góc $BOC = 90^\circ$ (hình vẽ); $\alpha$ là góc lượng giác hợp bởi tia đầu OA, tia cuối OB.
		Xét tính đúng sai các mệnh đề sau:}{\includegraphics[scale=.5]{img/HXN-4.13}}
	\choiceTF
	{\True Chiều cao của B so với mặt đất là $h_B = 23 + 20\sin \alpha$ ($mét$)}
	{\True Khi $\alpha = 45^\circ$ thì chiều cao của B so với mặt đất là $37,14$ $m$ (làm tròn kết quả đến hàng phần trăm)}
	{\True Chiều cao của C so với mặt đất là $h_C = 23 - 20\cos \alpha$ ($mét$)}
	{ Khi B ở vị trí có độ cao 33 $m$ thì C ở độ cao 13 $m$ so với mặt đất?}
	\loigiai{
	\begin{itemchoice}
		\itemch Chiều cao của B là $h_B = 23 + 20\sin \alpha$ ($mét$).
		\itemch Với $\alpha = 45^\circ$ thì $h_B = 23 + 20\sin 45^\circ \approx 37,14$ $m$.
		\itemch Chiều cao của C là $h_C = 23 + 20\sin (OA, OC) = 23 + 20\sin (\alpha - 90^\circ) = 23 - 20\cos \alpha$ ($mét$).
		\itemch Khi B ở vị trí có độ cao 33 $m$ thì $h_B = 23 + 20\sin \alpha = 33 \Rightarrow \sin \alpha = \dfrac{1}{2}$.\\
		Khi đó $\cos \alpha = \pm \sqrt{1-\sin^2 \alpha} = \pm \dfrac{\sqrt{3}}{2}$.\\
		Do đó $h_C = 23 - 20\cos \alpha =23\pm 10\sqrt{3}$.
	\end{itemchoice}
	}
\end{ex}
\begin{ex}
	Song sinh có thể là cùng trứng (identical) hoặc khác trứng (fraternal). Biết rằng $\dfrac{1}{3}$ số cặp song sinh là cùng trứng. Hiển nhiên, song sinh cùng trứng phải cùng giới tính; song sinh khác trứng có thể cùng hoặc khác giới tính. Giả sử song sinh cùng trứng có xác suất là hai bé trai hoặc hai bé gái như nhau, trong khi với song sinh khác trứng thì tất cả bốn khả năng đều có xác suất như nhau. Một nhà khảo sát tìm gặp ngẫu nhiên một người phụ nữ đang mang thai đôi.
	Xét tính đúng sai các mệnh đề sau:
	\choiceTF
	{\True Xác suất để người phụ nữ mang thai đôi là bé gái bằng $0,5$ biết rằng đây là cặp song sinh cùng trứng}
	{ Xác suất để thai đôi của người phụ nữ là một cặp trai gái bằng $0,3$}
	{Xác suất để thai đôi không cùng trứng và cũng không phải con trai bằng $\dfrac{1}{3}$}
	{\True Xác suất để người phụ nữ mang thai đôi là cùng trứng bằng $0,5$ biết rằng cô ấy hạ sinh được hai bé gái}
	\loigiai{
		Gọi A là biến cố: “Người phụ nữ mang thai đôi cùng trứng”, ký hiệu BB, GG, BG, GB lần lượt chỉ các biến cố thai đôi là trai-trai, gái-gái, trai-gái, gái trai. Ta có sơ đồ hình cây sau:
        \begin{center}
			\includegraphics[scale=.8]{img/HXN-4.14a}
		\end{center}
        \begin{itemchoice}
            \itemch Mệnh đề đúng. Ta có $P(GG|A) = \dfrac{1}{2}$.
			\itemch Mệnh đề sai. Ta có $P(BG \cup GB | \bar{A}) = \dfrac{2}{3} (\dfrac{1}{4} + \dfrac{1}{4}) = \dfrac{1}{3}$.
			\itemch Mệnh đề sai. Ta có $P(\bar{A} \cap GG) = \dfrac{2}{3} \cdot \dfrac{1}{4} = \dfrac{1}{6}$.
			\itemch Mệnh đề đúng. Ta có $P(GG) = \dfrac{1}{3} \cdot \dfrac{1}{2} + \dfrac{2}{3} \cdot \dfrac{1}{4}= \dfrac{1}{3}$.\\
			Do đó $P(A|GG) = \dfrac{P(A) \cdot P(GG|A)}{P(GG)} = \dfrac{\dfrac{1}{3} \cdot \dfrac{1}{2}}{\dfrac{1}{3}}= 0,5$.
        \end{itemchoice}
	}
\end{ex}
\begin{ex}
	\immini[thm]{ Một bể bơi hình trụ có đường kính $5$m và chiều cao $1$ m; bể được bơm nước vào với tốc độ không đổi $v_0$. Sau khi nước được bơm đầy, bể bị thủng một lỗ ở đáy và nước chảy ra ngoài; bể chảy hết nước trong $8$ giờ. Biết tốc độ giảm chiều cao của bể khi nước chảy ra ngoài vào thời điểm $t$ giờ (tính từ lúc nước đầy bể và ngừng bơm) được cho bởi hàm số $h'(t) = at + b$, với $a, b \in \mathbb{R}$. Lúc nước chảy hết ra ngoài thì tốc độ giảm chiều cao bằng $0$.
		Xét tính đúng sai các mệnh đề sau:}{\includegraphics[scale=.8]{img/HXN-4.15}}
	\choiceTF
	{\True Thể tích của bể bơi sau khi nước được làm đầy là $6,25\pi$ m$^3$}
	{$32a + 1 = 0$ và $4b - 1 = 0$}
	{\True Sau $4$ giờ kể từ lúc bể bị rò, lượng nước bị mất đi bằng $\dfrac{75\pi}{16}$ m$^3$}
	{\True Lượng nước bị rò rỉ ra ngoài một nửa sau $8 - 4\sqrt{2}$ giờ}
	\loigiai{
		Bể nước hình trụ có bán kính đáy $r = 2,5\ m$, chiều cao $h=1\ m$.
		\begin{itemchoice}
			\itemch Thể tích khi đầy là $V = \pi r^2 h = \pi \cdot (2,5)^2 \cdot 1 = 6,25\pi\ m^3$.
			\itemch Ta có $h'(8)=0 \Rightarrow 8a+b=0\ (1)$.\\ Chiều cao của nước thời điểm $t$ là $h(t) = \int\limits_{ }^{ }(at+b)\mathrm{d} t = \dfrac{at^2}{2} + bt + c$. \\Vì $h(0)=1 \Rightarrow c=1$. $h(8)=0 \Rightarrow 32a+8b+1=0\ (2)$. \\Từ (1) và (2) suy ra $a = \dfrac{1}{32}$, $b = -\dfrac{1}{4}$. Khi đó $32a-1=0$; $4b+1=0$ và $h(t) = \dfrac{1}{64}t^2 - \dfrac{1}{4}t + 1$.
			\itemch Chiều cao mực nước trong bể sau 4 giờ là: $h(4) = \dfrac{1}{64} \cdot 4^2 - \dfrac{1}{4} \cdot 4 + 1 = \dfrac{16}{64} - 1 + 1 = \dfrac{1}{4} = 0,25\ m$. Lượng nước còn lại trong bể sau 4 giờ là $\pi r^2 h(4) = \pi \cdot (2,5)^2 \cdot 0,25 = 6,25\pi \cdot 0,25 = 6,25\pi \cdot \dfrac{1}{4} = \dfrac{25\pi}{16}\ m^3$. Lượng nước đã thoát ra sau 4 giờ là $6,25\pi - \dfrac{25\pi}{16} = \dfrac{100\pi}{16} - \dfrac{25\pi}{16} = \dfrac{75\pi}{16}\ m^3$.
			\itemch Lượng nước còn lại khi bể mất một nửa nước là $\dfrac{6,25\pi}{2} = \dfrac{25\pi}{8}\ m^3$. \\Chiều cao tương ứng $h(t_1)$ của bể thỏa mãn $\pi r^2 h(t_1) = \dfrac{25\pi}{8} \Rightarrow h(t_1) = 0,5\ m$.\\ Ta có $h(t_1) = \dfrac{1}{64}t_1^2 - \dfrac{1}{4}t_1 + 1 = 0,5 \Leftrightarrow \dfrac{1}{64}t_1^2 - \dfrac{1}{4}t_1 + 0,5 = 0$. $\left[ \begin{array}{l} t_1 = 8 + 4\sqrt{2} \approx 13,7 > 8 \\ t_1 = 8 - 4\sqrt{2} \approx 2,3 \in (0; 8) \end{array} \right.$. \\Ta thấy $t_1 = 8 - 4\sqrt{2}$ (giờ) thỏa mãn đề bài.
		\end{itemchoice}
	}
\end{ex}
\begin{ex}
	Trong một mô hình game 3D, với hệ trục tọa độ thích hợp, người chơi đứng cùng với khẩu súng của anh ta được mô phỏng như một chất điểm di chuyển trên mặt phẳng $(P): x - 2y + 2z - 3 = 0$ và nhắm bắn các mục tiêu di động trên mặt cầu $(S)$ có phương trình $x^2 + y^2 + z^2 +2x - 4y - 2z + 5 = 0$. Người chơi vẫn có thể bắn trúng mục tiêu nếu nó di chuyển trên bán cầu khuất phía sau tầm nhìn. Sau khi trò chơi bắt đầu, anh ta quyết định nhắm bắn theo phương vector $\vec{u} = (1; 0; 1)$.
	Xét tính đúng sai các mệnh đề sau:
	\choiceTF
	{\True Mặt cầu $(S)$ có tâm $I(-1; 2; 1)$ và bán kính $R = 1$}
	{\True Mặt phẳng $(P)$ và mặt cầu $(S)$ không có điểm chung}
	{\True Người chơi đứng ở vị trí giao điểm của $(P)$ và $Ox$, khoảng cách từ tâm quả cầu đến đường bay viên đạn bằng $\dfrac{\sqrt{66}}{2}$}
	{\True Khoảng cách lớn nhất từ vị trí người bắn đến mục tiêu bằng $3\sqrt{2}$}
	\loigiai{
		\begin{enumerate}
			\itemch Mệnh đề đúng. Mặt cầu $(S)$ có tâm $I(-1; 2; 1)$ và bán kính $R = 1$.
			\itemch Mệnh đề đúng. Mặt phẳng $(P)$ và mặt cầu $(S)$ không có điểm chung.\\
			      Ta có $d(I, (P)) = \dfrac{|-1-4+2-3|}{\sqrt{1+4+4}} = 2$. $2 \ge R$. Do đó $(P)$ và mặt cầu $(S)$ không có điểm chung.
			\itemch Mệnh đề đúng. \\
			      Giao điểm của $(P)$ và $Ox$ là điểm $M(0; 3; 0) \Rightarrow \vec{IM} = (4; -2; -1)$; $[\vec{IM}, \vec{u}] = (-2; -5; 2)$.\\
			      Khoảng cách từ $I$ đến đường bay viên đạn là $d = \dfrac{|[\vec{IM}, \vec{u}]|}{|\vec{u}|} = \dfrac{\sqrt{(-2)^2 + (-5)^2 + 2^2}}{\sqrt{1^2 + 0^2 + 1^2}} =\dfrac{\sqrt{66}}{2}$.
			\itemch Mệnh đề đúng. Gọi $M$ thuộc $(P)$, $N$ thuộc $(S)$ theo thứ tự là vị trí người chơi và vị trí mục tiêu đang bắn; $H$ là hình chiếu của điểm $N$ trên $(P)$.\\
			      \begin{center}
				      \includegraphics[scale=1]{img/HXN-4.16a}
			      \end{center}
			      $MN$ hợp với $(P)$ một góc $\phi$ thỏa mãn $\sin \phi = \dfrac{|\vec{u}\cdot\vec{n_P}|}{|\vec{u}|\cdot|\vec{n_P}|} = \dfrac{|1+0+2|}{\sqrt{1+0+1}\cdot\sqrt{1+4+4}} = \dfrac{\sqrt{2}}{2}$.\\
			      Xét tam giác $MNH$ vuông tại $H$, ta có
			      $\sin \varphi = \dfrac{NH}{MN} \Rightarrow MN = \dfrac{NH}{\sin \varphi}$ hay $\boxed{MN = \sqrt{2}NH}$.
			      Dễ thấy $MN$ lớn nhất khi và chỉ khi $NH$ lớn nhất; mà
			      $NH \le d(I, (P)) + R = 2+1=3$.
			      Do đó $MN$ lớn nhất bằng $\boxed{3\sqrt{2}}$; khi đó $N, I, H$ nằm trên đường thẳng vuông góc với $(P)$.
		\end{enumerate}
	}
\end{ex}

\Closesolutionfile{ans}
\caukq
\Opensolutionfile{ans}[ans/ans-HXN-\sode-SA]


\begin{ex}
	\immini[thm]{ Một tấm cầu dốc kê bậc thềm được làm bằng kim loại như hình vẽ. Biết chiều cao tối đa của cầu dốc là $0,3 \, m$ và bề mặt cầu là hình vuông có cạnh bằng $1 \, m$. Hãy tính góc tạo bởi đường chéo bề mặt cầu dốc với mặt phẳng sàn nhà theo đơn vị độ (làm tròn kết quả đến hàng phần chục).
		\shortans{12,2}}{\includegraphics[scale=.8]{img/HXN-4.17}}
	\loigiai{
		Xét mô hình cầu dốc với các kí hiệu như hình vẽ.
		Vì $AF$ là hình chiếu của $BF$ trên mặt phẳng $(ACFD)$ nên góc giữa $(BF, (ACFD))$ là góc $\angle BFA$.
		Hình vuông $BCFE$ có cạnh bằng $1 \, m$ nên đường chéo $BF = \sqrt{1^2 + 1^2} = \sqrt{2} \, m$; $AB = 0,3 \, m$.
		Tam giác $ABF$ vuông tại $A$ có:
		$$ \sin \angle BFA = \frac{AB}{BF} = \frac{0,3}{\sqrt{2}} = \frac{0,3 \sqrt{2}}{2} = \frac{3 \sqrt{2}}{20} $$
		$$ \Rightarrow \angle BFA \approx 12,2^\circ $$
		Vậy góc giữa $(BF, (ACFD)) = \angle BFA \approx 12,2^\circ$.
	}
\end{ex}
\begin{ex}%Câu 1
	Có hai vợ chồng đã nghĩ ra một trò chơi đầy trí tuệ như sau: Họ sử dụng hai ly nước giống hệt nhau, mỗi ly chứa tối đa 240 ml nước. Ban đầu, người vợ có một ly nước đầy và người chồng có một cái ly rỗng. Bước thứ nhất người vợ rót 1/2 lượng nước trong ly của mình sang ly của người chồng; bước tiếp theo người chồng lại rót 1/3 lượng nước trong ly của mình sang cho ly người vợ. Quá trình này cứ tiếp tục mà mỗi lần rót thì mẫu số được cộng thêm 1; trò chơi này hấp dẫn đến mức cả hai người thực hiện đến bước thứ 100 thì dừng lại, hỏi lượng nước trong ly người chồng khi đó là bao nhiêu ml? (Làm tròn kết quả đến hàng đơn vị, giả sử trong quá trình rót nước không có giọt nước nào tràn ra ngoài).
	\shortans{119}
	\loigiai{
		Ta có thể thực hiện việc rót theo sơ đồ sau:
		\begin{center}
			\includegraphics[scale=.5]{img/HXN-4.18a}
		\end{center}
		Quá trình này được lặp đi lặp lại và ta thấy rằng trong các bước lẻ (người vợ rót nước cho người chồng) thì lượng nước hai ly bằng nhau.\\
		• Bước thứ 99 thì lượng nước hai ly bằng nhau.\\
		• Bước thứ 100 (người chồng rót 1/101 nước trong ly cho vợ), lượng nước trong ly người chồng là $\dfrac{1}{2}V-\dfrac{1}{101}\cdot\dfrac{1}{2}V=\dfrac{50V}{101}=\dfrac{50\cdot 240}{101}\approx\,119\,ml$ .}
\end{ex}
\begin{ex}%Câu 2
	Cho tập hợp $X=\left\{ 3\,;\,\,4\,;\,\,5\,;\,\,6\right\}$ và Y là tập hợp tất cả số tự nhiên có 2025 chữ số lấy từ X. Chọn ngẫu nhiên một số trong tập Y, biết rằng xác suất để số đó chia hết cho 3 bằng $\dfrac{1}{3}\left(\dfrac{1}{2^a}+b\right)$ , trong đó $a\,,\,\,b$ là các số nguyên dương. Tính giá trị $a-18b$ .
	\shortans{4031}
	\loigiai{
	Gọi $A_n, B_n$ là tập con của $Y$ gồm các số có $n$ chữ số với $A_n$ là tập các số chia hết cho 3 và $B_n$ là tập các số không chia hết cho 3.
	\begin{itemize}
		\item Với mỗi số thuộc $A_n$, có hai cách thêm vào cuối một chữ số 3 hoặc 6 để được $A_{n+1}$ và hai cách thêm vào cuối chữ số 4 hoặc chữ số 5 để được $B_{n+1}$.
		      \begin{itemize}
			      \item Ví dụ: $36 \in A_2$, nếu ta thêm 3 hoặc 6 vào sau nó thì được 363 hoặc 366 đều thuộc $A_3$.
			      \item Ví dụ: $36 \in A_2$, nếu ta thêm 4 hoặc 5 vào sau nó thì được 364 hoặc 365 đều thuộc $B_3$.
		      \end{itemize}
		\item Với mỗi số thuộc $B_n$, có một cách thêm vào cuối một chữ số 4 (hoặc chữ số 5) để được $A_{n+1}$ và có ba cách thêm một số để được $B_{n+1}$.
		      \begin{itemize}
			      \item Ví dụ: $34 \in B_2$, nếu ta thêm 5 vào sau nó thì được 345 thuộc $A_3$.
			      \item Ví dụ: $34 \in B_2$, nếu ta thêm 3 hoặc 4 hoặc 6 vào sau nó thì được 343 hoặc 344 hoặc 346 đều thuộc $B_3$.
		      \end{itemize}
	\end{itemize}
	Do đó ta có:
	\begin{align} |A_{n+1}| &= 2|A_n| + |B_n| & \\ |B_{n+1}| &= 2|A_n| + 3|B_n| & \end{align}
	Thay (1) vào (2), ta được: $|B_{n+1}| = 2|A_n| + 3(|A_{n+1}| - 2|A_n|) = 3|A_{n+1}| - 4|A_n| \quad (3)$.
	\\ Từ (3) suy ra $|B_n| = 3|A_n| - 4|A_{n-1}| \quad (4)$.
	\\Thay (4) vào (1) suy ra $|A_{n+1}| = 2|A_n| + 3|A_n| - 4|A_{n-1}| = 5|A_n| - 4|A_{n-1}|$.
	\\ Do đó $|A_{n+1}| = 5|A_n| - 4|A_{n-1}|$. Xét phương trình đặc trưng $t^2 = 5t - 4 \Rightarrow t=1 \lor t=4$.
	\\Phương trình $(*)$ có nghiệm dạng $|A_n| = a \cdot 1^n + b \cdot 4^n$. $|A_n| = a + b \cdot 4^n$ (a, b là các tham số).
	$A_1 = \{3; 6\} \Rightarrow |A_1| = 2 \Rightarrow a+4b = 2 \quad (5)$.
	\\$A_2 = \{36; 63; 33; 66; 45; 54\} \Rightarrow |A_2| = 6 \Rightarrow a+16b = 6 \quad (6)$.
		\\Từ (5) và (6) ta giải ra được $a = \dfrac{2}{3}$, $b = \dfrac{1}{3}$.
		\\Do đó $|A_{2025}| = \dfrac{2+4^{2025}}{3}$.
		\\Xác suất số đó chia hết cho 3 từ tập $Y$ là:
	$\dfrac{2+4^{2025}}{3} \times \dfrac{1}{4^{2025}} = \dfrac{1}{3} (\dfrac{2}{3 \cdot 4^{2050}} + \dfrac{1}{3 \cdot 4^{2051}})$.
		\\ Do đó $a = 4049, b = 1 \Rightarrow |A_{n}| = a - 18b = 4031$.}
\end{ex}
\begin{ex}
	\immini[thm]{ Cho hàm số $y = -\dfrac{1}{3}x^3 + \dfrac{3}{4}x^2 + 3x$ có đồ thị $(C)$ và đường thẳng $d$ đi qua gốc tọa độ tạo thành hai miền phẳng có diện tích $S_1$ và $S_2$ như hình vẽ. Biết $S_1 = \dfrac{27}{4}$ và $S_2 = \dfrac{m}{n}$ (hai số $m, n$ là nguyên tố cùng nhau), tính giá trị $2m-n$.\shortans{142}}{\includegraphics[scale=.8]{img/HXN-4.20}}
	\loigiai{
		\begin{center}
			\includegraphics[scale=.7]{img/HXN-4.20a}
		\end{center}
		Gọi $a>0$ là hoành độ giao điểm của $(C)$ và $d$.\\
		Đường thẳng $d$ có hệ số góc $k =\dfrac{-\dfrac{1}{3}a^3 + \dfrac{3}{4}a + 3a}{a} = -\dfrac{1}{2}a^2+\dfrac{3}{4}a + 3$.\\
		Mặt khác $d$ đi qua gốc tọa độ nên phương trình là $y = \left(-\dfrac{1}{2}a^2+\dfrac{3}{4}a + 3\right)x$.\\
		Ta có: $S_1 = \int\limits_{0}^{a} \left[-\dfrac{1}{2}x^3 +\dfrac{3}{4}x^2 + 3x - \left(-\dfrac{1}{2}a^2 + \dfrac{3}{4}a + 3\right)x\right] \mathrm{d} x$
		$\Leftrightarrow \dfrac{27}{4} = \dfrac{1}{8}a^4-\dfrac{1}{8}a^3 \Rightarrow a=3 > 0$.
		Ta có phương trình $d: y = \dfrac{3}{4}x$.
		Khi đó phương trình hoành độ giao điểm của $(C)$ và $d$ là:
		$\dfrac{3}{4}x - \left(-\dfrac{1}{2}x^3 + \dfrac{3}{4}x^2 + 3x\right) = 0 \Leftrightarrow x= 3 ; x = 0 ; x = -\dfrac{3}{2}$.\\
		Do đó $S_2 = \int\limits_{-\frac{3}{2}}^{0} \left|\dfrac{1}{2}x^3 - \dfrac{3}{4}x^2 - \dfrac{9}{4}x\right| \mathrm{d} x = \dfrac{135}{128} = \dfrac{m}{n}$. Từ đó suy ra $2m-n=142$.
	}
\end{ex}
\begin{ex}%Câu 4
	\immini[thm]{Một mảnh đất hình chữ nhật có kích thước $40\,m\times 50\,m$ đang được người chủ trồng cỏ tự nhiên. Vào buổi sáng, khi mặt trời vừa lên, mảnh đất này bị một mái nhà xưởng gần đó chắn ánh sáng. Khi mặt trời lên cao hơn, ánh sáng đã chiếu từ từ lên mảnh đất. Ta xem ranh giới giữa phần được chiếu sáng và phần tối là các đường thẳng song song thay đổi.
		Có thời điểm đường ranh giới này đi qua hai điểm A, M như hình vẽ (M là trung điểm một cạnh hình chữ nhật).
		Khi diện tích phần tối của mảnh đất bằng $75\,\,m^2$ , người ta đo được tốc độ giảm cạnh theo phương AD bằng 2 cm/s; hỏi tốc độ giảm diện tích phần tối của mảnh đất là bao nhiêu $c{m^2}$ /s? Kết quả được làm tròn đến hàng phần chục.\shortans{3873}}{\includegraphics[scale=1.1]{img/HXN-4.21}}

	\loigiai{
	\begin{center}
		\includegraphics[scale=1]{img/HXN-4.21a}
	\end{center}
	Xét tam giác ABM vuông tại B có $\tan\widehat{BAM}=\dfrac{BM}{AB}=\dfrac{2}{5}$ .\\
	Đặt $x=DF\in\left[0\,;\,\,40\right]\,,\,\,y=DE\in\left[0\,;\,\,50\right]$ (x, y thay đổi (giảm) vì ánh sáng ngày càng lan rộng).\\
	Vì $EF\text{//}AM$ nên $\tan\widehat{FED}=\tan\widehat{BAM}\Leftrightarrow\dfrac{x}{y}=\dfrac{2}{5}\Leftrightarrow y=\dfrac{5x}{2}$ .\\
	Diện tích phần tối tức thời là $S=\dfrac{1}{2}xy=\dfrac{1}{2}x\cdot\dfrac{5x}{2}=\dfrac{5}{4}x^2$ (1).\\
	Khi diện tích phần tối bằng $75\,\,m^2$ thì $\dfrac{5}{4}{x^2}=75\Rightarrow{x^2}=60\Rightarrow x=2\sqrt{15}$ m.\\
	Đạo hàm hai vế của (1) theo biến t ta được: $ \dfrac{dS}{dt}=\dfrac{5}{2}x .\dfrac{dx}{dt}\quad (2)$ .\\
	Thay $x=2\sqrt{15}\,\,m=200\sqrt{15}\,\,cm$ và $\dfrac{dx}{dt}=2$ cm/s vào (2), ta được $\dfrac{dS}{dt}\approx 3873\,\,c{m^2}$ /s.}
\end{ex}
\begin{ex}%Câu 5
	Trong không gian $Oxyz$ , cho điểm $A\left(-2\,;\,\,6\,;\,\,0\right)$ và mặt phẳng $\left(\alpha\right):3x+4y+89=0$ .
	Đường thẳng $d$ thay đổi nằm trên mặt phẳng $\left(Oxy\right)$ và luôn đi qua điểm $A$ .
	Gọi $H$ là hình chiếu vuông góc của $M\left(4\,;\,\,-2\,;\,\,3\right)$ trên đường thẳng d.
	Khoảng cách nhỏ nhất từ $H$ đến mặt phẳng $\left(\alpha\right)$ bằng bao nhiêu?
	\shortans{15}
	\loigiai{
	\begin{center}
		\includegraphics[scale=1]{img/HXN-4.22}
	\end{center}
	Gọi $K$ là hình chiếu vuông góc của $M$ lên $\left(Oxy\right)$ , suy ra $K\left(4\,;\,\,-2\,;\,\,0\right)$ .\\
	Vì $\left\{\begin{aligned}
			 & AH\perp MK\,\,\,\left(\text{do}\,\,\,MK\perp\left(Oxy\right)\right) \\
			 & AH\perp MH                                                          \\
		\end{aligned}\right.\Rightarrow AH\perp\left(MKH\right)\Rightarrow AH \perp KH $ .\\
	Khi đó $H$ luôn thuộc đường tròn $(C)$ có tâm là trung điểm $I\left(1\,;\,\,2\,;\,\,0\right)$ của đoạn $AK$ , bán kính $R=\dfrac{AK}{2}=5$ .\\
	Gọi $\Delta=\left(\alpha\right)\cap\left(Oxy\right)$ , ta thấy $\left(\alpha\right)\perp\left(Oxy\right)$ (vì $\vec{n}_{\left(\alpha\right)}\cdot\vec{k}=0$) . Khi đó: $d\left(I\,,\,\,\Delta\right)=d\left(I\,,\,\,\left(\alpha\right)\right)=20$ .\\
	Khoảng cách nhỏ nhất từ $H$ đến mặt phẳng $\left(\alpha\right)$ là $d{\left(H\,,\,\,\left(\alpha\right)\right)_{\text{min}}}=d{\left(H\,,\,\,\Delta\right)_{\text{min}}}=d\left(H_0\,,\,\,\Delta\right)$ $=d\left(I\,,\,\,\Delta\right)-R=20-5=15$ .}
\end{ex}
\Closesolutionfile{ans}
\inputansbox{6,4,3}{ans/ans-HXN-\sode-T,ans/ans-HXN-\sode-TF,ans/ans-HXN-\sode-SA}